


\section{Eftersøgning}

\begin{frame}
  \vspace{3em}
  \begin{center}
    {\LARGE \textbf{Bekendtgørelse om eftersøgning og aflivning af nødstedt vildt}} \\
    \vspace{1em}
    BEK nr 827 af 24/06/2016
  \end{center}
\end{frame}

\begin{frame}{Eftersøgning af anskudt klovvildt – pligt til handling (§1)}
    \begin{itemize}
        \item \textbf{§1 stk. 1}: Ved tegn på anskydning af klovbærende vildt under jagt
          \begin{itemize}
            \item Hvis vildtet ikke er fundet inden \textbf{6 dagtimer} (solopgang til solnedgang) efter anskydning
            \item Skal indehaver(e) af jagtretten tilkalde en \textbf{bemyndiget hundefører} (fra Naturstyrelsen) med henblik på aflivning
          \end{itemize}
        \item \textbf{Stk. 2}: Hvis indehaver(e) ikke selv deltog i jagten → pligten påhviler \textbf{den, der anskød vildtet}
        \item \textbf{Formål}: Sikre hurtig aflivning og minimere dyrets lidelser (jf. jagtetik)
    \end{itemize}
\end{frame}

\begin{frame}{Bemyndigelse af hundeførere og rettigheder (§2–§3)}
    \begin{itemize}
        \item \textbf{§2 stk. 1}: Naturstyrelsen bemyndiger hundeførere (og andre) til at eftersøge og aflive
          \begin{itemize}
            \item Anskudt eller på anden måde nødstedt vildt (inkl. olieindsmurte fugle)
          \end{itemize}
        \item \textbf{Stk. 2}: Bemyndigede har adgang til privat område uden retskendelse (mod forevisning af legitimation) – jf. lovens §11 stk. 1
        \item \textbf{Stk. 3}: Under eftersøgning efter §1: Kun den godkendte \textbf{schweisshund} må medtages
          \begin{itemize}
            \item Hundefører må medtage jagtvåben + \textbf{én ubevæbnet hjælper}
          \end{itemize}
        \item \textbf{§3}: Naturstyrelsen udsteder \textbf{legitimationskort} – skal forevises på forlangende
        \item \textbf{Schweissudvalg}: Rådgiver Naturstyrelsen om bemyndigelser og retningslinjer
    \end{itemize}
\end{frame}

\begin{frame}{Praktiske procedurer ved eftersøgning (§4)}
    \begin{itemize}
        \item \textbf{§4 stk. 1}: Inden eftersøgning starter i område, hvor hundefører/rekvirent \textbf{ikke} er jagtberettiget
          \begin{itemize}
            \item Skal hundeføreren \textbf{underrette politiet} om eftersøgningen
          \end{itemize}
        \item \textbf{Stk. 2}: Hvis vildtet ikke findes → underret grundejer/jagtberettiget om forløbet
        \item \textbf{Stk. 3}: Hvis vildtet findes → underret grundejer/jagtberettiget om placering (så de kan overtage vildtet)
        \item Husk: Schweissregistret (via DJ) bruges ofte til at finde nærmeste bemyndiget hundefører
    \end{itemize}
\end{frame}

\begin{frame}{Aflivning og retningslinjer (§5–§6)}
    \begin{itemize}
        \item \textbf{§5 stk. 1}: Aflivning skal ske \textbf{så smertefrit og hurtigt som muligt}
        \item \textbf{Stk. 2}: Legitimerede hundeførere må bruge \textbf{glatløbet haglgevær} til aflivning – også hvor hagl normalt er forbudt for arten
        \item \textbf{Stk. 3}: Må bruge \textbf{ikke-ekspanderende} riffelammunition – også hvor ekspanderende normalt kræves
        \item \textbf{§6}: Naturstyrelsen fastsætter nærmere \textbf{retningslinjer} for eftersøgning og aflivning
          \begin{itemize}
            \item Inkl. schweissarbejde, udstyr og procedurer (findes på Naturstyrelsens hjemmeside)
          \end{itemize}
    \end{itemize}
\end{frame}